% Gesamttest «Elektrizität» — Quelle des PDFs (python3 scripts/build-lp-pdf.py elektrizitaet).
% Musterlösung und Raster: bewertungspaket.tex. Alle Werte mit python3 nachgerechnet (03.10.2026).
\documentclass[11pt]{article}
\usepackage{../lp-druck}
\renewcommand{\lpname}{Elektrizität}
\renewcommand{\lpteilgebiet}{6.2}
\renewcommand{\lpfuss}{Gesamttest · 25 Punkte}
\begin{document}
\lpkopf{Gesamttest}{%
  \vspace{4pt}\noindent\begin{tabularx}{\linewidth}{@{}X X@{}}
  Code (kein Name): \hrulefill & Punkte: \hrulefill\ / 25
  \end{tabularx}}

\begin{lpinfo}{Anleitung}
\textbf{Zeit:} rund 25 Minuten. \textbf{Hilfsmittel:} Taschenrechner und Formelsammlung, in allen Teilen.
Schreib zu jeder Aufgabe zuerst die Formel, dann die Werte \textbf{mit Einheiten} — bewertet wird der Weg.
Danach: Lösung scannen oder fotografieren (ohne Namen und Standort) und mit dem \emph{Bewertungspaket} bewerten lassen.
\end{lpinfo}

\lpteil{Teil A · Ladung, Strom, Energie}{Kapitel 1 und 2 · 12 Punkte}

\begin{lpaufgabe}{G1}{Ein geriebener Ballon}{3 P}
Ein Ballon wird an den Haaren gerieben und trägt danach die Ladung $Q = -8.0\;\text{nC}$ ($e = 1.602 \cdot 10^{-19}\;\text{C}$).
Hat er Elektronen aufgenommen oder abgegeben, und wie viele? Welche Ladung tragen danach die Haare?
\lpplatz{5}
\end{lpaufgabe}

\begin{lpaufgabe}{G2}{E-Bike-Akku}{4 P}
Ein E-Bike-Akku ist mit $36\;\text{V}$ und $14\;\text{Ah}$ beschriftet.
Welcher Ladung in Coulomb entsprechen $14\;\text{Ah}$? Wie lange liefert der Akku $7\;\text{A}$?
Welche Energie in Wattstunden speichert er ($W = U \cdot Q$)?
\lpplatz{7}
\end{lpaufgabe}

\begin{lpaufgabe}{G3}{Backofen}{5 P}
Ein Backofen am Netz ($230\;\text{V}$) zieht $9.5\;\text{A}$. Berechne seine Leistung.
Er läuft $45\;\text{min}$: Wie viel Energie setzt er um, in Kilowattstunden, und was kostet das bei $0.27\;\text{CHF/kWh}$?
Jemand sagt: «Der Backofen verbraucht $9.5\;\text{A}$.» Was ist daran falsch?
\lpplatz{8}
\end{lpaufgabe}

\lpteil{Teil B · Widerstand und Schaltungen}{Kapitel 3 und 4 · 10 Punkte}

\begin{lpaufgabe}{G4}{Verlängerungskabel}{4 P}
Ein Verlängerungskabel aus Kupfer ($\rho = 0.017\;\Omega\,\text{mm}^2/\text{m}$) ist $20\;\text{m}$ lang und hat zwei Adern zu je $1.5\;\text{mm}^2$.
Welchen Widerstand hat die Leitung? Ein Heizlüfter zieht $8\;\text{A}$ durch das Kabel: Welche Spannung fällt an der Leitung ab?
\lpplatz{7}
\end{lpaufgabe}

\begin{lpaufgabe}{G5}{Zwei Widerstände}{6 P}
$R_1 = 60\;\Omega$ und $R_2 = 120\;\Omega$ liegen an $U = 24\;\text{V}$.
\begin{enumerate}[label=\alph*),leftmargin=*,itemsep=2pt,topsep=2pt]
\item In Reihe: Berechne $R_\text{ges}$, die Stromstärke $I$ und die Teilspannungen $U_1$ und $U_2$.
\item Parallel: Berechne $R_\text{ges}$, die Teilströme $I_1$ und $I_2$ und den Gesamtstrom $I$.
\end{enumerate}
\lpplatz{10}
\end{lpaufgabe}

\lpteil{Teil C · Gefahren und Schutz}{Kapitel 5 · 3 Punkte}

\begin{lpaufgabe}{G6}{Defektes Gehäuse}{3 P}
Jemand fasst mit nassen Händen ($R_\text{K} = 1.2\;\text{k}\Omega$) an das defekte Metallgehäuse eines Geräts, dessen Schutzleiter fehlt; das Gehäuse liegt an $230\;\text{V}$.
Wie gross ist der Körperstrom? Was tun ein FI-Schutzschalter ($30\;\text{mA}$) und ein Leitungsschutzschalter ($13\;\text{A}$)? Begründe.
\lpplatz{6}
\end{lpaufgabe}

\vfill
{\small\color{lpgrau}Bewerten: mit dem Bewertungspaket (Musterlösung, Punkteraster, Auftrag an eine KI) — erst nach dem Lösen öffnen.}
\end{document}
